\chapter{Appendix B: A Reading List}

Fifty-five posts worth reading, curated from 5,233.

This list draws from the longest posts, the most cross-disciplinary, the signature topics, and the interviews. It is not comprehensive (nothing could be), but it provides entry points across the blog's seventeen-year range.

All posts are available at johndcook.com/blog.

\bigskip

\section*{Start Here: 10 Essential Posts}

New to the blog? Start with these ten posts, chosen to represent the range and character of seventeen years of writing.

\begin{longtable}{l p{4cm} p{3.5cm}}
\toprule
\textbf{Date} & \textbf{Title} & \textbf{Why It's Essential} \tabularnewline
\midrule
2013-09-24 & Interview with Sir Michael Atiyah & The longest post: a Fields Medalist on mathematical intuition \tabularnewline
2010-10-06 & Probability that a number is prime & Deep dive into prime distribution: signature topic, signature depth \tabularnewline
2009-10-02 & Circle of fifths and number theory & Math meets music; shows how everything becomes a math problem \tabularnewline
2018-09-04 & Pi primes & Playful exploration of the most-mentioned constant \tabularnewline
2023-03-17 & Fixed points of the Fourier transform & Technical but clear; exemplifies the "translation layer" \tabularnewline
2025-11-16 & Elliptic curve pairings in cryptography & The 2019 spike's legacy: number theory made urgent \tabularnewline
2016-02-20 & The empty middle: why no one is average & Statistics with insight; accessible to non-specialists \tabularnewline
2020-04-19 & LaTeX command frequencies & Meta-analysis of tooling; characteristic data-driven curiosity \tabularnewline
2011-04-12 & Evaluating weather forecast accuracy: an interview with Eric Floehr & Entrepreneurship from curiosity; shows range beyond pure math \tabularnewline
2008-01-09 & Moore's law and software bloat & The very first post: 182 words that started 17 years \tabularnewline
\bottomrule
\end{longtable}

\bigskip

\section*{The Interviews}

The blog's conversations with remarkable people. These are often the longest posts in the archive.

\begin{longtable}{l p{5.5cm} r}
\toprule
\textbf{Date} & \textbf{Title} & \textbf{Words} \tabularnewline
\midrule
2013-09-24 & Interview with Sir Michael Atiyah & 1,926 \tabularnewline
2010-09-13 & Applied topology and Dante: an interview with Robert Ghrist & 1,814 \tabularnewline
2009-08-30 & Interview with author Cliff Pickover & 1,594 \tabularnewline
2013-03-19 & Interview with Sacha Chua & 1,243 \tabularnewline
2011-04-12 & Evaluating weather forecast accuracy: an interview with Eric Floehr & 1,241 \tabularnewline
2012-07-30 & Interview with Michael Hammel, author of The Artist's Guide to GIMP & 1,100 \tabularnewline
2016-06-10 & Interview with Chris Toomey of Upcase & 1,095 \tabularnewline
2009-06-22 & Dan Bricklin interview & 1,028 \tabularnewline
2009-07-29 & Carl Franklin interview & 824 \tabularnewline
2010-03-08 & Interview with Clojure author & 425 \tabularnewline
\bottomrule
\end{longtable}

\bigskip

\section*{The Longest Posts}

Deep dives that grew beyond the typical length. Two of the top ten are interviews.

\begin{longtable}{l p{5.5cm} r}
\toprule
\textbf{Date} & \textbf{Title} & \textbf{Words} \tabularnewline
\midrule
2024-03-26 & Hallucinations of AI Science Models & 1,710 \tabularnewline
2022-09-21 & Field of order 9 & 1,695 \tabularnewline
2014-02-24 & Writing down an unwritten language & 1,601 \tabularnewline
2018-04-11 & Anatomy of a posit number & 1,575 \tabularnewline
2021-03-28 & Solving for Möbius transformation coefficients & 1,456 \tabularnewline
2018-05-11 & Two-sample t-test and robustness & 1,429 \tabularnewline
2010-10-06 & Probability that a number is prime & 1,390 \tabularnewline
\bottomrule
\end{longtable}

\bigskip

\section*{Pi: The Constant}

The most-mentioned number in the archive, explored from multiple angles.

\begin{longtable}{l p{5.5cm} r}
\toprule
\textbf{Date} & \textbf{Title} & \textbf{Words} \tabularnewline
\midrule
2018-09-04 & Pi primes & 548 \tabularnewline
2021-02-12 & Finding pi in pi with Perl & 488 \tabularnewline
2021-10-07 & Gamma and the Pi function & 451 \tabularnewline
2017-10-23 & Finding numbers in pi & 891 \tabularnewline
2011-03-14 & Algorithm used for world record pi calculations & 406 \tabularnewline
\bottomrule
\end{longtable}

\bigskip

\section*{Primes: The Building Blocks}

The integers that divide only by themselves and one, explored for seventeen years.

\begin{longtable}{l p{5.5cm} r}
\toprule
\textbf{Date} & \textbf{Title} & \textbf{Words} \tabularnewline
\midrule
2010-10-06 & Probability that a number is prime & 1,390 \tabularnewline
2025-08-11 & Memorable Primes & 434 \tabularnewline
2024-02-10 & Factoring pseudoprimes & 790 \tabularnewline
2023-09-26 & Analogy between prime numbers and simple groups & 576 \tabularnewline
2024-12-31 & Details of generating primes for cryptography & 410 \tabularnewline
\bottomrule
\end{longtable}

\bigskip

\section*{The Gamma Function}

The most important function not on a calculator.

\begin{longtable}{l p{5.5cm} r}
\toprule
\textbf{Date} & \textbf{Title} & \textbf{Words} \tabularnewline
\midrule
2025-01-16 & Dimensional analysis for gamma function values & 621 \tabularnewline
2023-01-07 & Maxwell-Boltzmann and Gamma & 635 \tabularnewline
2021-11-06 & Gamma distribution tail probability approximation & 686 \tabularnewline
2021-03-26 & Simple approximation for Gamma function & 449 \tabularnewline
\bottomrule
\end{longtable}

\bigskip

\section*{Fourier Analysis}

Decomposing signals into frequencies: the workhorse of modern engineering.

\begin{longtable}{l p{5.5cm} r}
\toprule
\textbf{Date} & \textbf{Title} & \textbf{Words} \tabularnewline
\midrule
2023-03-17 & Fixed points of the Fourier transform & 460 \tabularnewline
2022-04-13 & Hilbert transform and Fourier series & 645 \tabularnewline
2022-03-20 & Reverse engineering Fourier conventions & 780 \tabularnewline
2021-06-07 & Fourier, Gauss, and Heisenberg & 421 \tabularnewline
\bottomrule
\end{longtable}

\bigskip

\section*{Cryptography}

The 2019 spike and beyond: number theory made urgent.

\begin{longtable}{l p{5.5cm} r}
\toprule
\textbf{Date} & \textbf{Title} & \textbf{Words} \tabularnewline
\midrule
2025-11-16 & Elliptic curve pairings in cryptography & 671 \tabularnewline
2025-11-17 & Three-party Diffie-Hellman in one shot & 513 \tabularnewline
2025-09-30 & Time needed to factor large integers & 487 \tabularnewline
2025-10-06 & A quiet change to RSA & 449 \tabularnewline
2025-11-18 & RSA as a pairing & 423 \tabularnewline
\bottomrule
\end{longtable}

\bigskip

\section*{Statistics}

Probability and inference: the second-largest domain.

\begin{longtable}{l p{5.5cm} r}
\toprule
\textbf{Date} & \textbf{Title} & \textbf{Words} \tabularnewline
\midrule
2018-05-11 & Two-sample t-test and robustness & 1,429 \tabularnewline
2025-04-28 & Variance of variances. All is variance. & 614 \tabularnewline
2016-02-20 & The empty middle: why no one is average & 547 \tabularnewline
2009-11-24 & A beta-like distribution & 553 \tabularnewline
\bottomrule
\end{longtable}

\bigskip

\section*{Music}

Where mathematics meets sound.

\begin{longtable}{l p{5.5cm} r}
\toprule
\textbf{Date} & \textbf{Title} & \textbf{Words} \tabularnewline
\midrule
2009-10-02 & Circle of fifths and number theory & 706 \tabularnewline
2009-03-15 & Typesetting music with LilyPond & 612 \tabularnewline
\bottomrule
\end{longtable}

\bigskip

\section*{Computing and Programming}

The bridge between mathematics and machines.

\begin{longtable}{l p{5.5cm} r}
\toprule
\textbf{Date} & \textbf{Title} & \textbf{Words} \tabularnewline
\midrule
2020-04-19 & LaTeX command frequencies & 514 \tabularnewline
2021-10-02 & Tables and interpolation & 433 \tabularnewline
2009-01-23 & Probability distributions and object oriented programming & 947 \tabularnewline
\bottomrule
\end{longtable}

\bigskip

\section*{Cross-Disciplinary Posts}

Posts spanning multiple categories: the connections between fields.

\begin{longtable}{l p{4cm} p{3.5cm}}
\toprule
\textbf{Date} & \textbf{Title} & \textbf{Categories} \tabularnewline
\midrule
2025-08-27 & Why are CUDA kernels hard to optimize? & AI, Computing, HPC, ML, Performance \tabularnewline
2025-07-09 & Experiences with Nvidia & AI, Computing, HPC, Parallel, Software \tabularnewline
2025-06-05 & Scientific papers: innovation … or imitation? & AI, Creativity, Productivity, Science \tabularnewline
2008-11-17 & Fast way to test whether a number is a square & Computing, Math, Science, Software \tabularnewline
\bottomrule
\end{longtable}

\bigskip

\section*{Science and the Physical World}

Where mathematics touches reality.

\begin{longtable}{l p{5.5cm} r}
\toprule
\textbf{Date} & \textbf{Title} & \textbf{Words} \tabularnewline
\midrule
2017-04-09 & Mercury and the bandwagon effect & 476 \tabularnewline
2011-09-22 & Big data and humility & 353 \tabularnewline
2010-02-19 & Popular research areas produce more false results & 313 \tabularnewline
\bottomrule
\end{longtable}

\bigskip

\section*{How to Use This List}

This list represents roughly 1\% of the archive. It favors longer posts, which tend to be more developed, and posts that span categories, which tend to be more unusual.

For any specific topic---quaternions, continued fractions, Unicode, orbital mechanics---search the blog directly. The archive is searchable at johndcook.com/blog, and seventeen years of posts mean that most mathematical and computational topics have been touched at least once.

The interviews are particularly recommended as entry points. They reveal both the subjects' thinking and the interviewer's interests: what he chose to ask, where he pushed for detail, what connections he noticed.

For readers who want to follow the blog's evolution, reading chronologically from 2008 is illuminating. The early posts are shorter and more scattered; the recent posts are longer and more focused. Watching the transformation unfold gives a sense of intellectual development that no single post can provide.

\bigskip

\section*{A Note on Selection}

Any list of fifty-five posts from 5,233 involves judgment calls. These selections favor:

\begin{itemize}
\item Length (longer posts generally contain more developed arguments)
\item Cross-disciplinarity (posts spanning categories are rarer and often more interesting)
\item Signature topics (pi, primes, gamma, Fourier---the recurring obsessions)
\item Interviews (the only posts featuring other voices)
\end{itemize}

A different curator would choose different posts. The archive is deep enough to support many valid reading lists. This one provides starting points; the exploration is up to the reader.
